% Fragmento público generado desde propietarios íntegros.
% Residencia: 52.1.1.
% Fuente: ../incoming/radion_monografia_20260827/manuscrito/sections/02_genealogia.tex:1-180
\noindent La genealogía operatoria antecede a la medida angular y determina las variables conservadas por cada lector.\par\medskip

\subsubsection{Soporte APP--TRIT--TPK y construcción del estado enriquecido}

\paragraph{Antecedencia del estado enriquecido respecto de la publicación circular}

La definición escolar de radián comienza con un círculo ya dado, un centro ya
dado y una longitud de arco ya mensurable. Si el radio y el arco tienen la
misma longitud, el ángulo subtendido mide un radián. La definición es correcta
dentro de la geometría euclídea, pero no responde a una pregunta anterior:
¿qué estructura hace posible que una vuelta posea fase, orientación, acción y
memoria, y por qué la unidad natural se relaciona con \(2\pi\)?

HMT invierte el orden causal. Primero construye el estado finito, después su
prolongación coinductiva, luego la vuelta y sólo al final su publicación
angular. El radión no corrige el radián convencional; lo levanta a su dominio
genealógico.

\begin{typebox}
\[
\APP\to\TRIT\to\TPK\to X_{\mathrm{enr}}
\to\mathcal C^{\mathrm{disc}}_{\mathrm{cont}}
\to\bigl(\pih,e_{\HMT},\varphi_{\HMT},\alphah,A,C^*,\Delta_4,\hret\bigr)
\to\mathfrak r_\sigma.
\]
La matemática convencional sólo aparece después como lenguaje de realización:
\(S^1\), elipse simpléctica, métrica de Klein, espacio de Hilbert, circuito
LC, pantalla de Minkowski y ecuaciones de campo.
\end{typebox}

\paragraph{Soporte bivariado de APP: hojas, residuos, cocientes y acarreo}

La Plataforma de Aritmética Posicional (APP) conserva simultáneamente dos
realizaciones de los símbolos \(1,\dots,9\): una hoja aditiva y una hoja
multiplicativa. No se reduce un número a su valor publicado. Junto al valor
permanecen hoja, orientación, cociente, resto, acarreo, frontera, ruta y
supervivencia.

El censo básico tiene tres estratos que jamás deben identificarse por mera
coincidencia cardinal:
\[
104,976\longrightarrow468\;U_6\longrightarrow243\;w_6
\subset\F_3^6,\qquad |\F_3^6|=729.
\]
Las \(468\) emisiones decimales, las \(243\) palabras visibles y el ambiente
de \(729\) palabras son objetos distintos. El ambiente será decisivo para la
normalización \(20/81\).

En base cúbica \(B=1000\), la resta orientada de dos palabras de tríadas se
realiza de derecha a izquierda:
\begin{equation}
u_i=t_i-v_i+c_{i+1},\qquad
r_i=u_i\bmod1000,\qquad
c_i=\frac{u_i-r_i}{1000}.
\label{eq:subacarreo}
\end{equation}
Fijados los operandos y el acarreo remoto, cociente y resto son únicos. Éste es
el operador que, más adelante, publicará las cifras de \(\epsone\). No hay
ningún decimal elegido después de ver el radián.

La completación
\[
10^3=729+243+27+1
\]
no es una curiosidad decimal. Expresa la conservación de la unidad al pasar
del ambiente ternario a una cámara de mil posiciones. La diferencia
\(1000-729=271\) es la corona completante; la identidad correcta es
\(999=729+270\), no \(729+271\).

\paragraph{Orientación ternaria y clasificación de los 3 regímenes operatorios del TRIT}

El TRIT asigna a cada estado local uno de tres regímenes orientados. En la
convención temporal activa,
\[
\begin{aligned}
+1&\longleftrightarrow\text{retorno, memoria, pasado},\\
0&\longleftrightarrow\text{umbral, presente},\\
-1&\longleftrightarrow\text{apertura bidireccional, futuro}.
\end{aligned}
\]
La realización cuadrática distingue:
\begin{equation}
J^2=-I,\qquad N^2=0,\qquad R^2=I.
\label{eq:tres-regimenes}
\end{equation}
Así aparecen, respectivamente, la rotación elíptica, el umbral parabólico y
la separación hiperbólica. No se mezclan tres métricas. Se conservan tres
acciones tipadas sobre un antecedente común.

La raíz interna de \(-1\) no se importa suponiendo de antemano el plano
complejo. En el bloque APP, dos proyecciones \(P,Q\) con razón
\(r=3/4=90/120\) producen
\[
J_{\mathrm{comm}}=\frac{[P,Q]}{\sqrt{r(1-r)}},\qquad J_{\mathrm{comm}}^2=-I,
\]
mientras
\[
R_{\mathrm{comm}}=2P-I,\qquad R_{\mathrm{comm}}^2=I,\qquad
R_{\mathrm{comm}}J_{\mathrm{comm}}=-J_{\mathrm{comm}}R_{\mathrm{comm}}.
\]
La pareja genera la forma mínima de \(\mathrm{Cl}_{1,1}\simeq M_2(\R)\):
rotación y apertura emergen juntas, pero siguen siendo operaciones diferentes.

\paragraph{Composición dinámica del TPK: selección, transporte, memoria y retorno}

El Topological Phase Núcleo no es un nombre para una caja de objetos. Es la
composición que selecciona fase, transporta cardinalidad, levanta hojas,
registra memoria y ejecuta el retorno. La dependencia fundamental es
\[
U_t\longrightarrow\Gamma_9\longrightarrow K_{\mathrm{ph}}.
\]
La holonomía nonádica \(\Gamma_9\) actúa sobre el estado enriquecido; la
publicación \(K_{\mathrm{ph}}\) es posterior y conserva memoria reducida. La
identidad \(K_{\mathrm{ph}}^9|_r=E_+E_\times\) no convierte a \(K_{\mathrm{ph}}\) en
generador.

El certificado nonádico conserva la cadena completa
\[
w_6\to w_{12}\to w_{18}\to w_{24}\to w_{30}
\to R_{36}\to G_9,
\]
con \(396\) niveles, \(2376\) trits, \(43\) vueltas, \(387\) pares
\((K,K+9)\) por canal y cinco regiones de \(\pi\). La fase cumple
\[
g(k+9)=g(k),
\]
pero el estado completo no retorna:
\[
B_{k+9}\ne B_k.
\]
Ésta es la semilla matemática de la hélice: la proyección angular cierra; la
memoria avanza.

\paragraph{Extensión temporal no escindida \(C_{12}\to C_{108}\to C_9\)}

La memoria de doce pasos se organiza mediante
\begin{equation}
0\longrightarrow C_{12}\longrightarrow C_{108}
\longrightarrow C_9\longrightarrow0.
\label{eq:extension-108}
\end{equation}
El cociclo
\[
c(r,r')=\left[\left\lfloor\frac{r+r'}9\right\rfloor\right]_{12}
\]
mide el acarreo entre retorno de fase y avance de memoria. Nueve pasos cierran
la fase observable \(C_9\); no reinician el registro dodecafásico.

Conviene separar cuatro objetos:
\begin{center}
\begin{tabularx}{.95\textwidth}{>{\bfseries}l X}
\toprule
\(C_{12}\) & grupo cíclico de doce posiciones;\\
\(\Rstate\) & registro enriquecido \((E_m,\tau_m,\sigma_m,\kappa_m,\Lambda_m)_{m=1}^{12}\);\\
\(\Theta_{108}\) & publicación cruda en el calendario de 108;\\
\(\Gamma_{108}\) & historia u holonomía levantada.\\
\bottomrule
\end{tabularx}
\end{center}
Por ello, la expresión vaga \enquote{post-\(C_{12}\)} se sustituirá por
\enquote{posterior al registro dodecafásico \(\Rstate\), en su lectura
sexádica} cuando ése sea el mapa efectivo.

\paragraph{Emergencia correlativa de las 5 construcciones del continuo desde el estado único}

La estructura discreta del continuo no se obtiene ensamblando cinco teorías
independientes. Sus cinco aspectos internos ---espectral, solenoidal, calibre,
cohomológico y determinantal--- emergen correlativamente del mismo estado
APP--TRIT--TPK. Esta obra no vuelve a demostrar el tratado integral, pero
preserva el hecho causal que necesita el radión: la sección coinductiva que
genera \(\pi,e,\varphi\) y \(\alpha\) no vive en una recta real desnuda; vive
en un objeto con fibras, relaciones, orientación, acarreo, memoria, terminal y
lector.

\begin{statusbox}
\textbf{Falsador de genealogía.} Si una fórmula posterior recibe
\(180/\pi\), \(\epsR\), una masa o una tabla metrológica para elegir una fase,
un coeficiente o una rama, deja de ser una generación HMT-first. El contraste
externo puede reconocer una salida; no puede seleccionarla.
\end{statusbox}
